\section{问题重述}

\subsection{问题背景}

热风烘干通过调节烘房温度与湿度，促使药材内部热量传递和水分迁移。工艺参数不当会造成干燥时间延长、内部水分不均或品质波动。题目以长 $25\,\mathrm{cm}$、初始半径 $2\,\mathrm{cm}$ 的近似圆柱形药材为对象，要求利用附件中的环境和尺寸数据，描述其径向温度场与干基含水率场，并确定烘干时间。

\subsection{问题要求}

\begin{enumerate}
  \item \textbf{预热平衡阶段：}利用附件~1 给出的烘房温度和水分浓度变化，以及附录~2 的常物性参数，求 $0$--$1800\,\mathrm{s}$ 内不同径向位置的药材温度和干基含水率，并生成 \texttt{result1.xlsx}。
  \item \textbf{整个烘干过程：}采用附录~3 的含水率、温度相关经验公式，描述预热平衡与恒温干燥阶段的热湿演化，给出前 $3\,\mathrm{h}$ 的规定结果，并生成 \texttt{result2.xlsx}。
  \item \textbf{烘干时长判定：}以药材各处干基含水率均低于 $0.15\,\mathrm{kg/kg}$ 为终止要求，确定首次达标的时间，并生成 \texttt{result3.xlsx}。
  \item \textbf{考虑尺寸变化：}利用附件~2 的药材半径数据和附录~4 的经验公式，将固定区域模型扩展到收缩区域，重新确定烘干时长，并生成 \texttt{result4.xlsx}。
\end{enumerate}

四问的核心递进关系是：问题一建立固定半径下的基础时空求解框架；问题二引入变物性和阶段衔接；问题三在全过程解上增加终点事件判定；问题四进一步改变计算区域和材料参数口径。

\draftnote{附件数据与四个 Excel 模板尚未加入当前工作区。拿到附件后，应核对时间覆盖范围、列名、单位、缺失值以及 result2.xlsx 要求的实际结束时间。}

